% Part: first-order-logic
% Chapter: beyond
% Section: other-logics

\documentclass[../../../include/open-logic-section]{subfiles}

\begin{document}

\olfileid{fol}{byd}{oth}

\olsection{نور منطقونه}

لکه چې ښايي تر اوسه مو معلومه کړې وي، د نوي منطق جوړول ګران نۀ دي۔
تاسو هم خپله نحو جوړولے، د استنتاج يو نظام پنځولے، او ورته معنٰی پوهنه
ترتيبولے شئ۔ که غواړئ د !!{derivation} نظام مو د معنٰی پوهنې په نسبت
بشپړ وي، ښايي لږ ځيرکتيا ته اړ شئ؛ او د نړۍ د قانع کولو دپاره هم څه
هڅه پکار ده چې منطق مو په رښتيا په زړه پورې دے۔ خو په بدل کښې ئې د
خپل منطق د رياضيکي او محاسبوي خاصيتونو په سپړلو ډېر ساعتونه پاکه او
خوندوره بوختيا لرلے شئ۔

وروستيو لسيزو د صوري منطقونو رښتينې چاودنه ليدلې ده۔ فزي منطق د
مبهمو خاصيتونو په هکله د استدلال د مدل کولو دپاره جوړ شوے۔ احتمالي
منطق د ناڅرګندتيا په هکله استدلال مدل کوي۔ ډيفالټ منطقونه او
نايکنواخته منطقونه د ماتېدونکو استدلالي بڼو د مدل کولو دپاره جوړ شوي؛
يعنې داسې ``معقول'' استنتاجونه چې د نويو معلوماتو په راتلو وروسته
ردېدلے شي۔ معرفتي منطقونه د پوهې په هکله استدلال، سببي منطقونه د سبب
او مسبب د اړيکو په هکله استدلال، او آن ``تکليفي'' منطقونه د اخلاقي
مکلفيتونو په هکله استدلال مدل کوي۔ دا نظامونه چې د فلسفي، رياضيکي يا
محاسبوي انګېزې له مخې معرفي شوي وي، ښايي د رياضيکي منطق، فلسفي منطق،
مصنوعي ځيرکتيا، ادراکي پوهنې يا بلې څانګې تر سرليک لاندې مطالعه شي۔

لړ همداسې اوږدېږي او امکانونه بې‌پايه ښکاري۔ ښايي هېڅکله د لايبنېڅ
دې خوب ته و نۀ رسېږو چې ټول انساني عقل محاسبې ته راکم کړو---خو دا مو
له هڅې څخه نۀ شي منع کولے۔

\end{document}
